\ifthenelse{\equal{\lang}{ko}}{%
본 연구는 \textbf{AI 거버넌스 정보학(AI Governance Informatics)}을 향한 첫 걸음을 내디디며, 
전 세계적으로 확장 가능하고 신뢰할 수 있는 규제 모니터링을 위한 기초를 마련하였다. 
우리는 \textbf{LexSense}를 제안하였으며, 이는 실시간 탐지, 의미 분석, LLM 기반 보고를 통합한 통합 프레임워크이다. 
EU, 미국, 한국 데이터셋에서 91\%의 정확도를 달성하였고, 모니터링 업무량을 82\% 줄였으며, 
보고 시간을 70\% 단축하면서 지연은 3.5초 미만으로 유지하였다. 
사용자 연구를 통해 더 높은 신뢰성과 효율성이 확인되었다. 

공정성 인식 미세조정, 편향 감사, 프라이버시 보호 메커니즘을 내장함으로써, 
우리는 효율성 향상과 윤리적 안전장치를 조화시켰다. 
향후 연구에서는 다국어 커버리지를 확장하고, 개념 드리프트 하에서의 연속 학습을 고도화하며, 
경량 피드백 프로토콜을 설계할 예정이다. 
종합적으로, 본 연구는 배포 가능한 컴플라이언스 인프라를 제공하고, 
NLP, 법학, 책임 있는 AI의 교차점에서 AI 거버넌스를 데이터 기반의 구체적 학문 분야로 정립한다. 
 

}{%
% -------------------- English Version --------------------

This study establishes \textit{AI Governance Informatics} as a new research domain by introducing \textbf{LexSense}, an LLM-driven framework for regulatory monitoring and compliance. Through the integration of real-time detection, semantic data acquisition, and LLM-based reporting, LexSense demonstrated both high accuracy (91\%) and significant efficiency gains (82\% workload reduction, 70\% faster reporting) across EU, U.S., and Korean datasets. The release of open-source code, datasets, and Dockerized environments further enhances reproducibility and practical applicability.  

Beyond technical achievements, this work contributes a conceptual foundation that bridges natural language processing, information systems, and legal governance. By embedding fairness-aware tuning, bias audits, and privacy-preserving mechanisms, LexSense advances not only system efficiency but also responsible deployment.  

Looking forward, we envision extending this framework to low-resource jurisdictions, strengthening long-term robustness through automated drift detection, and deepening theoretical modeling of governance processes within organizations. In doing so, LexSense can evolve from a practical compliance system into a cornerstone framework for \textit{responsible AI governance} in information systems research.  

}
